\chapter{C++}\label{ch:iv3:cpp}

\begin{ivtrains}
Durée de vie, possession et comportement indéfini d'abord, puis écriture au tableau blanc des classes de la bibliothèque standard elle-même~: un pointeur intelligent, un vecteur, un allocateur. Les énoncés précisent avec quels outils on les compile lorsque la réponse en dépend, à savoir ici C++20 avec g++ 11.4, libstdc++ et x86-64 Linux. Apprendre séparément ce que le langage garantit et ce que ce compilateur fait effectivement.
\end{ivtrains}

\begin{question}[$\star$ \iqroles{QD, Dev, LLE} \iqfirm{toutes} \iqlevel{junior}]\label{iq:iv3:cpp:1}
Dans le \cref{lst:iv3:cp:ptrconst}, remplacez le commentaire, tour à tour, par chacune de ces lignes~: \texttt{*p = 3;}, \texttt{p = \&b;}, \texttt{*q = 3;}, \texttt{q = \&b;}, \texttt{*r = 3;}, \texttt{r = \&b;}. Lesquelles compilent~? Comment lisez-vous une telle déclaration~?
\omcode{code/interviews-3/03-cpp/cpp/ptr_const.cpp}{1}{8}{Trois pointeurs const.}\label{lst:iv3:cp:ptrconst}
\end{question}

\begin{question}[$\star$ \iqroles{QD, Dev, LLE} \iqfirm{toutes} \iqlevel{junior}]\label{iq:iv3:cpp:2}
Qu'affiche le \cref{lst:iv3:cp:storage}, dans l'ordre~? Quel objet n'est jamais détruit~? Nommez la durée de stockage de chacun des quatre objets, et dites où chacun réside sur x86-64 Linux.
\omcode{code/interviews-3/03-cpp/cpp/storage.cpp}{1}{22}{Quatre sortes de durée de vie.}\label{lst:iv3:cp:storage}
\end{question}

\begin{question}[$\star$ \iqroles{QD, Dev, LLE} \iqfirm{toutes} \iqlevel{junior}]\label{iq:iv3:cpp:3}
En C++, quelle est la différence entre \texttt{class} et \texttt{struct}~? Dans le \cref{lst:iv3:cp:classstruct}, remplacez le commentaire par chacune de ces expressions~: \texttt{(void)S\{\}.x;}, \texttt{(void)C\{\}.x;}, \texttt{(void)FromC\{\}.get();} et \texttt{(void)FromS\{\}.x;}. Lesquelles compilent~?
\omcode{code/interviews-3/03-cpp/cpp/class_struct.cpp}{1}{8}{Accès par défaut.}\label{lst:iv3:cp:classstruct}
\end{question}

\begin{question}[$\star$ \iqroles{QD, Dev, LLE} \iqfirm{toutes} \iqlevel{junior}]\label{iq:iv3:cpp:4}
Quelles sont les deux lignes affichées par le \cref{lst:iv3:cp:ref}~? Que fait \texttt{r = b}, et en quoi cela diffère-t-il de \texttt{p = \&b}~? Donnez deux autres différences entre une référence et un pointeur.
\omcode{code/interviews-3/03-cpp/cpp/ref_rebind.cpp}{1}{13}{Une référence et un pointeur.}\label{lst:iv3:cp:ref}
\end{question}

\begin{question}[$\star$ \iqroles{QD, Dev, LLE} \iqfirm{toutes} \iqlevel{junior}]\label{iq:iv3:cpp:5}
Avec g++ 11.4 sur x86-64, un \texttt{int} de 4 octets et des pointeurs de 8 octets, quels trois nombres affiche le \cref{lst:iv3:cp:arrayparam}~? Pourquoi le \texttt{[10]} du paramètre n'aide-t-il pas~? Comment passez-vous un tableau en conservant sa longueur~?
\omcode{code/interviews-3/03-cpp/cpp/array_param.cpp}{2}{12}{Un tableau en paramètre.}\label{lst:iv3:cp:arrayparam}
\end{question}

\begin{question}[$\star$ \iqroles{QD, Dev, LLE} \iqfirm{toutes} \iqlevel{junior}]\label{iq:iv3:cpp:6}
Avec g++ sur x86-64, quelles sont les trois lignes affichées par le \cref{lst:iv3:cp:promote}~? Quel est le type de \texttt{a + b}~? Y a-t-il là un comportement indéfini~?
\omcode{code/interviews-3/03-cpp/cpp/promote.cpp}{1}{11}{Petits entiers.}\label{lst:iv3:cp:promote}
\end{question}

\begin{question}[$\star$ \iqroles{QD, Dev, LLE} \iqfirm{toutes} \iqlevel{junior}]\label{iq:iv3:cpp:7}
Quelles sont les deux lignes affichées par le \cref{lst:iv3:cp:vectorinit}~? Énoncez la règle qui tranche entre les deux sens de chaque couple d'arguments.
\omcode{code/interviews-3/03-cpp/cpp/vector_init.cpp}{1}{8}{Parenthèses et accolades.}\label{lst:iv3:cp:vectorinit}
\end{question}

\begin{question}[$\star$ \iqroles{QD, Dev, LLE} \iqfirm{toutes} \iqlevel{junior}]\label{iq:iv3:cpp:8}
Qu'affiche le \cref{lst:iv3:cp:memberorder}~? La liste d'initialisation nomme d'abord \texttt{b}. Qu'est-ce qui détermine l'ordre, et quel avertissement vous aurait-il prévenu~? Pourquoi cela importe-t-il si l'initialiseur de \texttt{b} lit \texttt{a}~?
\omcode{code/interviews-3/03-cpp/cpp/member_order.cpp}{2}{19}{L'ordre des membres.}\label{lst:iv3:cp:memberorder}
\end{question}

\begin{question}[$\star$ \iqroles{QD, Dev, LLE} \iqfirm{toutes} \iqlevel{junior}]\label{iq:iv3:cpp:9}
Quelles sont les cinq lignes affichées par le \cref{lst:iv3:cp:overload}~? Quel est le type de \texttt{nullptr}~? Pourquoi l'a-t-on ajouté au langage, alors que 0 désignait déjà un pointeur nul~?
\omcode{code/interviews-3/03-cpp/cpp/overload_zero.cpp}{1}{16}{Zéro et nullptr.}\label{lst:iv3:cp:overload}
\end{question}

\begin{question}[$\star$ \iqroles{QD, Dev, LLE, QR} \iqfirm{toutes} \iqlevel{junior}]\label{iq:iv3:cpp:10}
Qu'affiche le \cref{lst:iv3:cp:removeonly}~? Pourquoi \texttt{std::remove} ne modifie-t-il pas la taille du vecteur~? Quelle est la façon, en une seule ligne et en C++20, d'effectuer les deux étapes~?
\omcode{code/interviews-3/03-cpp/cpp/remove_only.cpp}{1}{13}{Retirer sans effacer.}\label{lst:iv3:cp:removeonly}
\end{question}

\begin{question}[$\star$ \iqroles{QD, Dev, LLE, QR} \iqfirm{toutes} \iqlevel{junior}]\label{iq:iv3:cpp:11}
Quelles sont les deux lignes affichées par le \cref{lst:iv3:cp:accumulate}~? Le code compile sans avertissement sous \texttt{-Wall -Wextra}. Pourquoi le premier nombre est-il faux~?
\omcode{code/interviews-3/03-cpp/cpp/accumulate_int.cpp}{1}{9}{La valeur initiale d'une somme.}\label{lst:iv3:cp:accumulate}
\end{question}

\begin{question}[$\star$ \iqroles{QD, Dev, LLE} \iqfirm{toutes} \iqlevel{junior}]\label{iq:iv3:cpp:12}
Quelles sont les trois lignes affichées par le \cref{lst:iv3:cp:rangecopy}~? Que laisse de côté \texttt{auto} lorsqu'il déduit un type à partir d'un \texttt{const int\&}~? Quand une boucle for basée sur un intervalle doit-elle utiliser \texttt{const auto\&}, et quand \texttt{auto} seul est-il préférable~?
\omcode{code/interviews-3/03-cpp/cpp/range_copy.cpp}{1}{14}{Ce que déduit auto.}\label{lst:iv3:cp:rangecopy}
\end{question}

\begin{question}[$\star\star$ \iqroles{QD, Dev, LLE} \iqfirm{toutes} \iqlevel{junior, confirmé}]\label{iq:iv3:cpp:13}
Le \Cref{lst:iv3:cp:bufferbad} compile sans avertissement et semble souvent s'exécuter. Qu'est-ce qui ne va pas, et que signalerait AddressSanitizer en premier~? Énoncez la règle des cinq, puis écrivez cette classe pour la respecter. Pourriez-vous éviter d'écrire l'un de ces membres~?
\omcode{code/interviews-3/03-cpp/cpp/buffer_bad.cpp}{1}{14}{Une classe qui possède de la mémoire.}\label{lst:iv3:cp:bufferbad}
\end{question}

\begin{question}[$\star\star$ \iqroles{QD, Dev, LLE} \iqfirm{toutes} \iqlevel{junior, confirmé}]\label{iq:iv3:cpp:14}
Quelles sont les deux lignes affichées par le \cref{lst:iv3:cp:nomove} en C++20~? \texttt{Logged} ne diffère de \texttt{Plain} que par un destructeur vide, ajouté pour faciliter la journalisation. Quels membres spéciaux cela a-t-il supprimés~? Le programme reste-t-il correct, et qu'a-t-il perdu~? Quelle seule ligne le corrige~?
\omcode{code/interviews-3/03-cpp/cpp/no_move.cpp}{1}{26}{Un destructeur vide.}\label{lst:iv3:cp:nomove}
\end{question}

\begin{question}[$\star\star$ \iqroles{QD, Dev, LLE} \iqfirm{toutes} \iqlevel{junior, confirmé}]\label{iq:iv3:cpp:15}
Avec libstdc++, où un vecteur part vide et double sa capacité, quelles sont les deux lignes affichées par le \cref{lst:iv3:cp:vecnoexcept}~? Les deux types ne diffèrent que par \texttt{noexcept} sur le constructeur de déplacement. Pourquoi le vecteur effectue-t-il une copie dans un cas~? Quelle garantie protège-t-elle~?
\omcode{code/interviews-3/03-cpp/cpp/vec_noexcept.cpp}{1}{22}{Croissance et noexcept.}\label{lst:iv3:cp:vecnoexcept}
\end{question}

\begin{question}[$\star\star$ \iqroles{QD, Dev, LLE} \iqfirm{toutes} \iqlevel{junior}]\label{iq:iv3:cpp:16}
Avec libstdc++, quelles sont les cinq lignes affichées par le \cref{lst:iv3:cp:movedfrom}~? Lesquelles de ces valeurs la norme garantit-elle, et lesquelles ne sont que ce que fait cette bibliothèque~? Que fait \texttt{std::move} lui-même à l'exécution~? Que pouvez-vous faire d'une \texttt{std::string} déplacée~?
\omcode{code/interviews-3/03-cpp/cpp/moved_from.cpp}{1}{21}{Ce qui reste après un déplacement.}\label{lst:iv3:cp:movedfrom}
\end{question}

\begin{question}[$\star\star$ \iqroles{QD, Dev, LLE} \iqfirm{toutes} \iqlevel{junior, confirmé}]\label{iq:iv3:cpp:17}
Avec g++ 11.4 en C++20, quelles sont les trois lignes affichées par le \cref{lst:iv3:cp:elision}~? Pourquoi écrire \texttt{std::move} sur le retour est-il pire que de l'omettre~? Pourquoi \texttt{either} effectue-t-il une copie~? Et comment \texttt{make} peut-il renvoyer un type sans constructeur de copie ni de déplacement utilisable~?
\omcode{code/interviews-3/03-cpp/cpp/elision.cpp}{2}{35}{Renvoyer une valeur.}\label{lst:iv3:cp:elision}
\end{question}

\begin{question}[$\star\star$ \iqroles{QD, Dev, LLE} \iqfirm{toutes} \iqlevel{junior, confirmé}]\label{iq:iv3:cpp:18}
Quelles sont les cinq lignes affichées par le \cref{lst:iv3:cp:forwarding}~? Pourquoi \texttt{bad} effectue-t-il toujours une copie, même lorsqu'on lui transmet un objet temporaire~? Que fait \texttt{std::forward<T>} que \texttt{std::move} ne fait pas, et qu'est-ce qui ne va pas si vous utilisez \texttt{std::move} dans \texttt{good}~?
\omcode{code/interviews-3/03-cpp/cpp/forwarding.cpp}{1}{21}{Transmettre un argument.}\label{lst:iv3:cp:forwarding}
\end{question}

\begin{question}[$\star\star$ \iqroles{QD, Dev, LLE} \iqfirm{toutes} \iqlevel{confirmé}]\label{iq:iv3:cpp:19}
Quelles sont les trois lignes affichées par le \cref{lst:iv3:cp:collapse}~? Lesquels parmi \texttt{T\&\&} dans \texttt{f}, \texttt{std::vector<T>\&\&} dans \texttt{g} et \texttt{auto\&\&} sont des références universelles, et comment le reconnaissez-vous~? Vers quel type \texttt{T} est-il déduit dans \texttt{f(i)}, et comment \texttt{int\& \&\&} devient-il un type légal~? Remplacez le commentaire par \texttt{std::vector<int> l; g(l);}~: compile-t-il~?
\omcode{code/interviews-3/03-cpp/cpp/collapse.cpp}{1}{23}{Deux sens de la double esperluette.}\label{lst:iv3:cp:collapse}
\end{question}

\begin{question}[$\star\star$ \iqroles{QD, Dev, LLE} \iqfirm{toutes} \iqlevel{junior, confirmé}]\label{iq:iv3:cpp:20}
Le \Cref{lst:iv3:cp:svtemp} compile sans avertissement avec g++ 11.4, \texttt{-Wall -Wextra}. Qu'est-ce qui ne va pas~? Que signale AddressSanitizer~? Le bogue apparaîtrait-il avec une chaîne de 5 caractères~? Comment prenez-vous un \texttt{string\_view} en toute sécurité, et dans quel cas est-il un bon type de paramètre mais un mauvais membre~?
\omcode{code/interviews-3/03-cpp/cpp/sv_temp.cpp}{1}{10}{Une vue sur un temporaire.}\label{lst:iv3:cp:svtemp}
\end{question}

\begin{question}[$\star\star$ \iqroles{QD, Dev, LLE} \iqfirm{toutes} \iqlevel{junior}]\label{iq:iv3:cpp:21}
Quelles sont les deux lignes affichées par le \cref{lst:iv3:cp:svcstr}~? Un \texttt{string\_view} est transmis à une fonction C par \texttt{.data()}. Pourquoi est-ce un bogue, et pourquoi cela passe-t-il souvent ses tests~? Y a-t-il là un comportement indéfini~?
\omcode{code/interviews-3/03-cpp/cpp/sv_cstr.cpp}{1}{11}{Une vue et la fin d'une chaîne.}\label{lst:iv3:cp:svcstr}
\end{question}

\begin{question}[$\star\star$ \iqroles{QD, Dev, LLE} \iqfirm{toutes} \iqlevel{junior}]\label{iq:iv3:cpp:22}
Qu'affiche le \cref{lst:iv3:cp:cycle}~? Quelle ligne ne s'affiche jamais, et pourquoi~? Que signale un détecteur de fuites~? Corrigez-le à l'aide de \texttt{std::weak\_ptr}, et dites comment vous utilisez l'autre nœud une fois que son pointeur est faible.
\omcode{code/interviews-3/03-cpp/cpp/cycle.cpp}{1}{15}{Deux nœuds qui se retiennent mutuellement.\protect\hfil\protect\break}\label{lst:iv3:cp:cycle}
\end{question}

\begin{question}[$\star\star$ \iqroles{QD, Dev, LLE} \iqfirm{toutes} \iqlevel{junior}]\label{iq:iv3:cpp:23}
Écrivez une version simplifiée de \texttt{std::unique\_ptr<T>} en C++20, sans suppresseur personnalisé~: un constructeur à partir d'un pointeur, un destructeur, la construction et l'affectation par déplacement, aucune copie, ainsi que \texttt{get}, \texttt{release}, \texttt{reset} et les opérateurs. Que se passe-t-il lors d'une auto-affectation par déplacement~? Pourquoi \texttt{reset} enregistre-t-il le nouveau pointeur avant de supprimer l'ancien~? Quelle est la taille de votre classe~?
\end{question}

\begin{question}[$\star\star$ \iqroles{QD, Dev, LLE} \iqfirm{toutes} \iqlevel{junior, confirmé}]\label{iq:iv3:cpp:24}
Un type contient un mégaoctet, et le \cref{lst:iv3:cp:weak} remplace \texttt{operator delete} pour compter les libérations. Avec libstdc++, qu'affiche-t-il~? Quand le mégaoctet est-il rendu dans chaque cas, et pourquoi les deux diffèrent-ils~? Donnez un cas où il faut éviter \texttt{make\_shared}, et le cas qui le favorise habituellement.
\omcode{code/interviews-3/03-cpp/cpp/make_shared_weak.cpp}{6}{38}{Quand la mémoire est rendue.}\label{lst:iv3:cp:weak}
\end{question}

\begin{question}[$\star\star$ \iqroles{QD, Dev, LLE} \iqfirm{toutes} \iqlevel{junior}]\label{iq:iv3:cpp:25}
Qu'a de faux le \cref{lst:iv3:cp:twoowners}, et que montrent les deux compteurs~? Une méthode d'une classe veut renvoyer un \texttt{shared\_ptr} vers son propre objet. Pourquoi \texttt{std::shared\_ptr<Order>(this)} constitue-t-il le même bogue~? Quelle est la bonne façon de procéder, et que se passe-t-il si l'objet n'appartient pas déjà à un \texttt{shared\_ptr}~?
\omcode{code/interviews-3/03-cpp/cpp/two_owners.cpp}{1}{14}{Deux propriétaires d'un même pointeur brut.}\label{lst:iv3:cp:twoowners}
\end{question}

\begin{question}[$\star\star$ \iqroles{QD, Dev, LLE} \iqfirm{toutes} \iqlevel{junior}]\label{iq:iv3:cpp:26}
Le \Cref{lst:iv3:cp:vecref} peut afficher 9 ou 1, ou autre chose. Pourquoi~? Que signale AddressSanitizer, et sur quelle ligne~? Donnez deux corrections, dont l'une conserve la référence. Le même code sur un \texttt{std::deque} a-t-il ce bogue~?
\omcode{code/interviews-3/03-cpp/cpp/vec_ref.cpp}{1}{10}{Une référence dans un vecteur.}\label{lst:iv3:cp:vecref}
\end{question}

\begin{question}[$\star\star$ \iqroles{QD, Dev, LLE} \iqfirm{toutes} \iqlevel{junior, confirmé}]\label{iq:iv3:cpp:27}
Avec libstdc++, qu'affiche le \cref{lst:iv3:cp:vecbool}~? Pourquoi \texttt{auto} sur un \texttt{vector<int>} et sur un \texttt{vector<bool>} ne se comportent-ils pas de la même façon~? Le code \texttt{bool\& b = v[0];} compile-t-il~? Deux threads écrivent chacun un indice différent d'un \texttt{vector<bool>}~: s'agit-il d'une course aux données, et en est-ce une pour un \texttt{vector<char>}~?
\omcode{code/interviews-3/03-cpp/cpp/vec_bool.cpp}{1}{13}{Un vecteur qui n'est pas comme les autres.}\label{lst:iv3:cp:vecbool}
\end{question}

\begin{question}[$\star\star$ \iqroles{QD, Dev, LLE} \iqfirm{toutes} \iqlevel{junior, confirmé}]\label{iq:iv3:cpp:28}
Avec g++ 11.4 et libstdc++, quelles sont les trois lignes affichées par le \cref{lst:iv3:cp:stability}~? Quels résultats la norme garantit-elle~? Qu'un rehachage invalide-t-il dans un \texttt{std::unordered\_map}~? Vous conservez des pointeurs vers des ordres contenus dans un conteneur~: quels conteneurs le permettent~?
\omcode{code/interviews-3/03-cpp/cpp/stability.cpp}{1}{24}{Un élément reste-t-il en place~?}\label{lst:iv3:cp:stability}
\end{question}

\begin{question}[$\star\star$ \iqroles{QD, Dev, LLE} \iqfirm{toutes} \iqlevel{junior, confirmé}]\label{iq:iv3:cpp:29}
Quelles sont les trois lignes affichées par le \cref{lst:iv3:cp:emplace}~? Quand \texttt{emplace\_back} est-il réellement meilleur que \texttt{push\_back}~? Remplacez le commentaire par \texttt{rows.push\_back(1000000);}~: que se passe-t-il, et pourquoi est-ce préférable à ce que faisait \texttt{emplace\_back} sur la ligne précédente~?
\omcode{code/interviews-3/03-cpp/cpp/emplace.cpp}{1}{22}{Deux façons d'ajouter un élément.}\label{lst:iv3:cp:emplace}
\end{question}

\begin{question}[$\star\star$ \iqroles{QD, Dev, LLE, QR} \iqfirm{toutes} \iqlevel{junior, confirmé}]\label{iq:iv3:cpp:30}
Le \Cref{lst:iv3:cp:sortle} trie 100 valeurs égales avec un comparateur écrit \texttt{a <= b}. Avec g++ 11.4, que se passe-t-il sous AddressSanitizer~? Quelle règle le comparateur a-t-il enfreinte~? Pourquoi le même code peut-il réussir sur des entrées courtes ou distinctes~? Comment écririez-vous un comparateur sur deux champs~?
\omcode{code/interviews-3/03-cpp/cpp/sort_le.cpp}{1}{9}{Un comparateur avec un signe égal.}\label{lst:iv3:cp:sortle}
\end{question}

\begin{question}[$\star\star$ \iqroles{QD, Dev, LLE} \iqfirm{toutes} \iqlevel{junior, confirmé}]\label{iq:iv3:cpp:31}
Avec g++ 11.4, \texttt{-std=c++20 -O2}, x86-64, qu'affiche le \cref{lst:iv3:cp:overflow}~? La première fonction est une tentative courante de détecter un dépassement, pour $b > 0$. Pourquoi est-elle fausse, et que fait le compilateur avec elle~? Que modifient UBSan et \texttt{-fwrapv}~? Écrivez un test portable qui n'utilise aucune primitive du compilateur.
\omcode{code/interviews-3/03-cpp/cpp/overflow_check.cpp}{1}{15}{Vérifier le dépassement trop tard.}\label{lst:iv3:cp:overflow}
\end{question}

\begin{question}[$\star\star$ \iqroles{QD, Dev, LLE} \iqfirm{toutes} \iqlevel{junior, confirmé}]\label{iq:iv3:cpp:32}
Avec g++ 11.4, C++20, \texttt{-O2}, x86-64, exécuté sans argument, quelles sont les quatre lignes affichées par le \cref{lst:iv3:cp:shiftabs}~? Pour chaque ligne, la valeur est-elle garantie par la norme ou indéfinie~? Que donne le décalage, et pourquoi~? Lesquelles UBSan signalerait-il~?
\omcode{code/interviews-3/03-cpp/cpp/shift_abs.cpp}{1}{13}{Quatre cas limites sur les entiers.}\label{lst:iv3:cp:shiftabs}
\end{question}

\begin{question}[$\star\star$ \iqroles{QD, Dev, LLE} \iqfirm{toutes} \iqlevel{junior, confirmé}]\label{iq:iv3:cpp:33}
Dans le \cref{lst:iv3:cp:lifetime}, lesquelles des références \texttt{one}, \texttt{two} et \texttt{three} sont sûres à lire après leur déclaration~? Énoncez la règle. Pourquoi le passage du temporaire par une fonction qui renvoie son propre argument met-il la règle en défaut~? Que signalerait AddressSanitizer pour la dernière~?
\omcode{code/interviews-3/03-cpp/cpp/lifetime.cpp}{1}{15}{Lier un temporaire à une référence.}\label{lst:iv3:cp:lifetime}
\end{question}

\begin{question}[$\star\star$ \iqroles{QD, Dev, LLE} \iqfirm{toutes} \iqlevel{junior, confirmé}]\label{iq:iv3:cpp:34}
Quelles sont les quatre lignes affichées par le \cref{lst:iv3:cp:defaultargs}~? Pourquoi la première affiche-t-elle \texttt{Derived 1}~? Pourquoi \texttt{d.g(1)} convertit-il 1 en \texttt{double} alors que \texttt{Base} offre une meilleure correspondance~? Donnez la correction pour chacun, et dites lequel peut être signalé par une option du compilateur.
\omcode{code/interviews-3/03-cpp/cpp/default_args.cpp}{1}{21}{Valeurs par défaut et masquage.}\label{lst:iv3:cp:defaultargs}
\end{question}

\begin{question}[$\star\star$ \iqroles{QD, Dev, LLE} \iqfirm{toutes} \iqlevel{junior, confirmé}]\label{iq:iv3:cpp:35}
Avec g++ 11.4, libstdc++ et x86-64, quelles sont les trois lignes affichées par le \cref{lst:iv3:cp:sizes}~? Pourquoi une classe vide n'a-t-elle pas une taille nulle, et pourquoi \texttt{Base1} occupe-t-elle moins de place que \texttt{Member}~? Pourquoi un suppresseur lambda ne coûte-t-il rien, alors qu'un suppresseur à pointeur de fonction coûte 8~octets~? Que représentent les 16 octets de \texttt{shared\_ptr}~?
\omcode{code/interviews-3/03-cpp/cpp/sizes.cpp}{1}{23}{Ce que les objets occupent.}\label{lst:iv3:cp:sizes}
\end{question}

\begin{question}[$\star\star$ \iqroles{QD, Dev, LLE} \iqfirm{toutes} \iqlevel{junior}]\label{iq:iv3:cpp:36}
Un en-tête déclare \texttt{template <class T> T total(const std::vector<T>\& v);}. Son corps se trouve dans \texttt{sum.cpp}, et \texttt{main.cpp} inclut l'en-tête et appelle \texttt{total} sur un \texttt{std::vector<int>}. Chaque fichier compile, puis la construction échoue. À quelle étape, et pourquoi~? Donnez les deux corrections. Que coûte la seconde à un utilisateur qui veut un \texttt{std::vector<long>}~?
\end{question}

\begin{question}[$\star\star$ \iqroles{QD, Dev, LLE} \iqfirm{toutes} \iqlevel{junior, confirmé}]\label{iq:iv3:cpp:37}
Quelles sont les deux lignes affichées par le \cref{lst:iv3:cp:concepts}~? Comment le compilateur choisit-il entre les deux versions de \texttt{kind}~? Remplacez le commentaire par \texttt{notional(1, 2);}~: que se passe-t-il, et où pointe l'erreur~? Comment auriez-vous écrit cette contrainte avant C++20~?
\omcode{code/interviews-3/03-cpp/cpp/concepts.cpp}{1}{25}{Contraindre un modèle.}\label{lst:iv3:cp:concepts}
\end{question}

\begin{question}[$\star\star$ \iqroles{QD, Dev, LLE} \iqfirm{crypto, teneur de marché, hedge fund} \iqlevel{junior, confirmé}]\label{iq:iv3:cpp:38}
Les messages de données de marché appartiennent à l'un de trois types, et le \cref{lst:iv3:cp:variant} les conserve dans un \texttt{std::variant}. Avec g++ 11.4 sur x86-64, quelles sont les quatre lignes qu'il affiche~? Comment cela se compare-t-il à une classe de base et trois classes dérivées, pour l'allocation et pour l'ajout d'un nouveau type~? Si le constructeur d'une nouvelle alternative lève une exception pendant \texttt{emplace}, un variant peut-il finir sans rien contenir~?
\omcode{code/interviews-3/03-cpp/cpp/msg_variant.cpp}{1}{31}{Des messages dans un variant.}\label{lst:iv3:cp:variant}
\end{question}

\begin{question}[$\star\star$ \iqroles{QD, Dev, LLE} \iqfirm{toutes} \iqlevel{junior, confirmé}]\label{iq:iv3:cpp:39}
Vous voulez un \texttt{std::unordered\_map} indexé par une structure formée d'un symbole et d'une plateforme de négociation. Quelles sont les deux choses que vous devez fournir, et que se passe-t-il si vous n'en fournissez aucune des deux~? Le \Cref{lst:iv3:cp:keyhash} le remplit avec 100 symboles sur 50 plateformes chacun, avec deux fonctions de hachage. Avec libstdc++, quel est le plus grand compartiment sous le hachage faible, et quelle en est à peu près la taille sous le hachage mélangé~? Le hachage faible donne-t-il de mauvais résultats~? Quel en est le coût~?
\omcode{code/interviews-3/03-cpp/cpp/key_hash.cpp}{1}{36}{Un hachage pour une structure.}\label{lst:iv3:cp:keyhash}
\end{question}

\begin{question}[$\star\star$ \iqroles{QD, Dev} \iqfirm{toutes} \iqlevel{junior}]\label{iq:iv3:cpp:40}
Les types valeur, vus depuis un autre langage. Avec .NET 8, quelles sont les quatre lignes affichées par le \cref{lst:iv3:cp:boxing}~? \texttt{Fill} est une \texttt{struct} C\#. Qu'est-ce que le boxing, quand a-t-il lieu ici, et pourquoi l'appel de \texttt{Add} à travers l'interface ne modifie-t-il pas \texttt{a}~? Le \Cref{lst:iv3:cp:valuecpp} est la version C++, qui affiche \texttt{3 1 9}~: que propose C++ à la place du boxing~?
\omcode{code/interviews-3/03-cpp/csharp/Boxing/Program.cs}{1}{24}{Une structure en C\#.}\label{lst:iv3:cp:boxing}
\omcode{code/interviews-3/03-cpp/cpp/value_cpp.cpp}{1}{15}{La même chose en C++.}\label{lst:iv3:cp:valuecpp}
\end{question}


\begin{question}[$\star\star\star$ \iqroles{QD, Dev, LLE} \iqfirm{toutes} \iqlevel{junior, confirmé}]\label{iq:iv3:cpp:41}
\texttt{Item} possède un constructeur de copie et un constructeur de déplacement \texttt{noexcept}, et chacun incrémente \texttt{Item::moves} d'une unité. Avec libstdc++, quelles sont les deux lignes affichées par le \cref{lst:iv3:cp:growth}~? Pourquoi un vecteur croît-il par un facteur, alors que croître de 16 éléments à chaque fois gaspillerait moins de mémoire~? Donnez le nombre total de déplacements pour 1000 insertions avec chaque politique. Pourquoi une bibliothèque pourrait-elle choisir un facteur inférieur à 2~? Puis le \cref{lst:iv3:cp:naivevec}~: que fait \texttt{v.push\_back(v[0])}, et comment le vrai vecteur doit-il l'éviter~?
\omcode{code/interviews-3/03-cpp/cpp/growth.cpp}{11}{22}{Observer la capacité.\protect\hfil\protect\break}\label{lst:iv3:cp:growth}
\omcode{code/interviews-3/03-cpp/cpp/naive_vec.cpp}{6}{37}{Une première tentative de vecteur.}\label{lst:iv3:cp:naivevec}
\end{question}

\begin{question}[$\star\star\star$ \iqroles{QD, Dev, LLE} \iqfirm{toutes} \iqlevel{confirmé}]\label{iq:iv3:cpp:42}
Écrivez une version simplifiée de \texttt{std::optional<T>} en C++20, avec \texttt{emplace}, \texttt{reset}, une conversion vers \texttt{bool}, \texttt{operator*}, ainsi que la copie et le déplacement. \texttt{T} n'a pas besoin de constructeur par défaut. Pourquoi ne peut-il pas simplement avoir un membre \texttt{T value}~? À quoi doit ressembler le stockage~? À quel moment faut-il activer le drapeau, et pourquoi~? Le \Cref{lst:iv3:cp:misaligned} place un \texttt{double} à une adresse impaire. x86-64 le charge sans broncher, le programme est-il donc correct~?
\omcode{code/interviews-3/03-cpp/cpp/misaligned.cpp}{4}{8}{Un objet à la mauvaise adresse.}\label{lst:iv3:cp:misaligned}
\end{question}

\begin{question}[$\star\star\star$ \iqroles{QD, Dev, LLE} \iqfirm{crypto, teneur de marché, hedge fund} \iqlevel{confirmé}]\label{iq:iv3:cpp:43}
Un système de trading crée et détruit de nombreux ordres d'un même type, et ne doit pas toucher au tas une fois en marche. En C++20, écrivez un pool à taille fixe \texttt{Pool<T, N>}~: \texttt{create(args.{}.{}.)} renvoie un \texttt{T*}, ou \texttt{nullptr} lorsqu'il est plein, et \texttt{destroy(p)} le reprend, les deux en temps constant. Où conservez-vous la liste libre des emplacements~? Que doit-il se passer si le constructeur de \texttt{T} lève une exception~? Quelles erreurs votre pool ne détecte-t-il pas~?
\end{question}

\begin{question}[$\star\star\star$ \iqroles{QD, Dev, LLE} \iqfirm{toutes} \iqlevel{confirmé}]\label{iq:iv3:cpp:44}
En C++20, écrivez une bump arena au-dessus d'un tampon de taille fixe, ainsi qu'un type d'allocateur qui permet à \texttt{std::vector<int>} et à \texttt{std::map<int, double>} de prendre toute leur mémoire dans celle-ci. Que doit fournir au minimum un allocateur~? Pourquoi a-t-il besoin d'un constructeur de conversion à partir d'un allocateur d'un autre type, alors que votre code ne nomme jamais que \texttt{pair}~? Que fait \texttt{deallocate}, et que deviennent les anciens blocs lorsqu'un vecteur grandit~?
\end{question}

\begin{question}[$\star\star\star$ \iqroles{QD, Dev, LLE} \iqfirm{toutes} \iqlevel{confirmé}]\label{iq:iv3:cpp:45}
Avec g++ 11.4 et libstdc++, quelles trois lignes affiche le \cref{lst:iv3:cp:pmr}~? Le programme remplace l'\texttt{operator new} global de façon à ajouter 1 à \texttt{news}. Expliquez 1020 et 400. Pourquoi les chaînes du second bloc utilisent-elles le tampon, alors que personne ne le leur a transmis~? Qu'apporte cette conception par rapport à un allocateur passé en argument de patron, et que coûte-t-elle~?
\omcode{code/interviews-3/03-cpp/cpp/pmr_use.cpp}{19}{57}{Ressources mémoire polymorphes. Ici \texttt{pmr} est un alias de \texttt{std::pmr}.\protect\hfil\protect\break}\label{lst:iv3:cp:pmr}
\end{question}

\begin{question}[$\star\star\star$ \iqroles{QD, Dev, LLE} \iqfirm{toutes} \iqlevel{junior, confirmé}]\label{iq:iv3:cpp:46}
Dans le \cref{lst:iv3:cp:assignbad}, \texttt{get} représente une allocation qui peut lever une exception. Trouvez deux défauts distincts dans \texttt{operator=}. Pour chacun, dites précisément ce qui se passe~: \texttt{a = a}, puis \texttt{a = b} lorsque \texttt{get} lève une exception. Réécrivez ensuite l'opérateur de sorte que les deux cas soient corrects. Que coûte votre version, et quand pourriez-vous choisir de ne pas la payer~?
\omcode{code/interviews-3/03-cpp/cpp/assign_bad.cpp}{1}{27}{Une affectation avec deux défauts.}\label{lst:iv3:cp:assignbad}
\end{question}

\begin{question}[$\star\star\star$ \iqroles{QD, Dev, LLE} \iqfirm{crypto, teneur de marché, hedge fund} \iqlevel{confirmé}]\label{iq:iv3:cpp:47}
Avec g++ 11.4, \texttt{-std=c++20 -O2}, x86-64, quelles deux lignes affiche le \cref{lst:iv3:cp:aliasing}~? Pourquoi la première ligne n'est-elle pas 0, alors que l'écriture par \texttt{b} a bien écrit des zéros sur \texttt{x}~? Que modifie \texttt{-fno-strict-aliasing}~? Vous analysez un prix de 4 octets dans un paquet, stocké dans un tampon \texttt{const char*}. Comment le lisez-vous correctement, et quel en est le coût~?
\omcode{code/interviews-3/03-cpp/cpp/aliasing.cpp}{1}{15}{Une adresse, deux types.}\label{lst:iv3:cp:aliasing}
\end{question}

\begin{question}[$\star\star\star$ \iqroles{QD, Dev, LLE} \iqfirm{toutes} \iqlevel{confirmé}]\label{iq:iv3:cpp:48}
Le \Cref{lst:iv3:cp:nullcheck} comporte deux fonctions qui ne diffèrent que par l'ordre de deux lignes. Avec g++ 11.4, \texttt{-O2}, x86-64, en quel code chacune se compile-t-elle~? Pourquoi le contrôle de nullité a-t-il été supprimé dans l'une d'elles~? Que se passe-t-il avec \texttt{read\_bad(nullptr)}~? Quelle option conserve le contrôle, et est-ce la bonne correction~?
\omcode{code/interviews-3/03-cpp/cpp/nullcheck.cpp}{1}{10}{Un contrôle placé après l'utilisation.}\label{lst:iv3:cp:nullcheck}
\end{question}

\begin{question}[$\star\star\star$ \iqroles{QD, Dev, LLE} \iqfirm{toutes} \iqlevel{junior, confirmé}]\label{iq:iv3:cpp:49}
Le \Cref{lst:iv3:cp:maxdangle} se compile sans avertissement avec g++ 11.4, \texttt{-Wall -Wextra}. Exécutez-le sans argument, puis avec trois. Sous AddressSanitizer, que se passe-t-il à chaque fois, et pourquoi le nombre d'arguments compte-t-il~? Pourquoi \texttt{std::max} renvoie-t-il une référence~? Pourquoi \texttt{std::max(1, 2.0)} ne compile-t-il pas~? Donnez la correction de ce listing.
\omcode{code/interviews-3/03-cpp/cpp/max_dangle.cpp}{1}{7}{Le plus grand des deux.}\label{lst:iv3:cp:maxdangle}
\end{question}

\begin{question}[$\star\star\star$ \iqroles{QD, Dev, LLE} \iqfirm{crypto, teneur de marché, hedge fund} \iqlevel{confirmé}]\label{iq:iv3:cpp:50}
Une configuration en service est conservée dans un \texttt{std::shared\_ptr}~: un thread la remplace, et un autre la copie pour la lire. Le compteur de références de \texttt{shared\_ptr} est atomique. Dans le \cref{lst:iv3:cp:sharedrace}, exécuté avec l'argument \texttt{A} sous ThreadSanitizer, le programme est-il exempt de course aux données~? Que peut-il se passer réellement~? Que change le mode \texttt{B}~? Quatre threads détiennent chacun leur propre copie d'un même \texttt{shared\_ptr} et la copient dans une boucle~: s'agit-il d'une course aux données~? Avec g++ 11.4, peut-on utiliser \texttt{std::atomic<std::shared\_ptr<int>{}>}~?
\omcode{code/interviews-3/03-cpp/cpp/shared_race.cpp}{1}{25}{Remplacer un pointeur partagé pendant la lecture.}\label{lst:iv3:cp:sharedrace}
\end{question}

\begin{question}[$\star\star\star$ \iqroles{QD, Dev, LLE} \iqfirm{toutes} \iqlevel{junior, confirmé}]\label{iq:iv3:cpp:51}
En C++20, écrivez une file bornée bloquante pour plusieurs producteurs et consommateurs, avec \texttt{push}, \texttt{pop} et \texttt{close}. \texttt{push} se bloque lorsque la file est pleine et renvoie \texttt{false} une fois fermée. \texttt{pop} se bloque lorsque la file est vide et renvoie un \texttt{std::optional} vide une fois fermée et vidée. Pourquoi deux variables de condition~? Pourquoi chaque attente prend-elle un prédicat~? Que doit faire \texttt{close} pour que personne ne dorme indéfiniment~? Comment le testeriez-vous~?
\end{question}

\begin{question}[$\star\star\star$ \iqroles{QD, Dev, LLE, QR} \iqfirm{toutes} \iqlevel{junior, confirmé}]\label{iq:iv3:cpp:52}
En C++20, écrivez \texttt{choose(n, k)} pour le coefficient binomial, sur \texttt{std::uint64\_t}, renvoyant un \texttt{std::optional} vide exactement lorsque la vraie réponse ne tient pas sur 64 bits. Il ne doit pas signaler de débordement pour une réponse qui tient, même si un produit intermédiaire naïf déborderait. Pourquoi la formule factorielle échoue-t-elle, et à partir de quel $n$~? Démontrez que vos valeurs intermédiaires ne débordent jamais, sauf si la réponse déborde. Quel est le plus petit $n$ pour lequel la réponse peut déborder~?
\end{question}

\begin{question}[$\star\star\star$ \iqroles{QD, Dev, LLE} \iqfirm{toutes} \iqlevel{junior, confirmé}]\label{iq:iv3:cpp:53}
Le \Cref{lst:iv3:cp:terminate} est exécuté trois fois sous g++ 11.4, avec pour argument \texttt{noexcept}, puis \texttt{destructor}, puis \texttt{plain}. Que se passe-t-il à chaque fois~? Pourquoi le bloc \texttt{try} ne rattrape-t-il pas les deux premiers~? Le destructeur de \texttt{Closer} ne porte pas de \texttt{noexcept}, alors pourquoi se comporte-t-il comme le premier~? Pourriez-vous laisser un destructeur lever une exception, et dans quel cas cela terminerait-il quand même le programme~?
\omcode{code/interviews-3/03-cpp/cpp/terminate.cpp}{2}{20}{Trois façons de lever une exception.}\label{lst:iv3:cp:terminate}
\end{question}

\begin{question}[$\star\star\star$ \iqroles{QD, Dev, LLE} \iqfirm{crypto, teneur de marché, hedge fund} \iqlevel{junior, confirmé}]\label{iq:iv3:cpp:54}
Qu'affiche le \cref{lst:iv3:cp:crtp}~? Expliquez comment \texttt{Strat<Momo>} atteint \texttt{Momo::signal} sans fonction virtuelle. Avec g++ 11.4, \texttt{-O2}, qu'est-ce que l'assembleur de \texttt{static\_call} contient de différent de celui de \texttt{virtual\_call}~? D'où vient réellement la rapidité~? Que se passe-t-il si \texttt{Momo} n'a pas de \texttt{signal}~? Pourquoi ne peut-on pas mettre un \texttt{Momo} et une autre stratégie de cette forme dans un même conteneur~?
\omcode{code/interviews-3/03-cpp/cpp/crtp.cpp}{1}{31}{Répartition statique et dynamique.}\label{lst:iv3:cp:crtp}
\end{question}

\begin{question}[$\star\star\star$ \iqroles{QD, Dev, LLE} \iqfirm{toutes} \iqlevel{junior, confirmé}]\label{iq:iv3:cpp:55}
Exécuté sans argument, qu'affiche le \cref{lst:iv3:cp:constexpr}~? Remplacez le commentaire successivement par \texttt{int d = twice(argc);}, \texttt{int e = always(argc);}, \texttt{constexpr int f = twice(argc);} et \texttt{constexpr int o = twice(2000000000);}. Lesquelles compilent~? En quoi diffèrent \texttt{constexpr} et \texttt{consteval} appliqués à une fonction~? La dernière ligne aurait un comportement indéfini à l'exécution, alors pourquoi s'agit-il ici d'une erreur de compilation~?
\omcode{code/interviews-3/03-cpp/cpp/constexpr_table.cpp}{1}{27}{Calculs à la compilation.}\label{lst:iv3:cp:constexpr}
\end{question}

\begin{question}[$\star\star\star\star$ \iqroles{QD, Dev, LLE} \iqfirm{toutes} \iqlevel{confirmé}]\label{iq:iv3:cpp:56}
En C++20, implémentez un \texttt{Vec<T>} avec le destructeur, la copie et le déplacement, \texttt{emplace\_back}, \texttt{pop\_back} et \texttt{operator[]}, et qui fonctionne pour un type sans constructeur par défaut. \texttt{emplace\_back} doit offrir la garantie forte, et \texttt{v.emplace\_back(v[0])} doit fonctionner lorsqu'il déclenche une réallocation. Pourquoi pas \texttt{new T[c]}~? Dans quel ordre s'enchaînent les étapes d'une réallocation~? Quand déplacez-vous les anciens éléments, et quand les copiez-vous~? Comment testeriez-vous la garantie~?
\end{question}

\begin{question}[$\star\star\star\star$ \iqroles{QD, Dev, LLE} \iqfirm{crypto, teneur de marché, hedge fund} \iqlevel{confirmé}]\label{iq:iv3:cpp:57}
En C++20, implémentez un pointeur à comptage de références \texttt{Rc<T>}, un \texttt{shared\_ptr} sans pointeurs faibles, avec une seule allocation pour le compteur et l'objet, et un compteur sûr lorsque plusieurs threads détiennent des copies. Quel ordre mémoire l'incrément requiert-il~? Lequel requiert la décrémentation, et que casse-t-on si elle est relaxed~? Montrez un test qui détecterait cela sur un x86-64 à ordre fort. Qu'est-ce que votre classe ne rend toujours pas sûr~?
\end{question}

\begin{question}[$\star\star\star\star$ \iqroles{QD, Dev, LLE} \iqfirm{crypto, teneur de marché, hedge fund} \iqlevel{confirmé}]\label{iq:iv3:cpp:58}
En C++20, implémentez \texttt{Fn<R(A.{}.{}.)>}~: un objet appelable à déplacement seul et à effacement de type, avec un petit tampon de 32 octets, qui se rabat sur le tas pour un objet plus grand ou dont le déplacement peut lever une exception. Il doit contenir une lambda qui capture un \texttt{std::unique\_ptr}. Comment appelez-vous une fonction dont vous avez oublié le type~? De quoi avez-vous besoin en plus de l'appel~? Pourquoi le tampon ne doit-il servir qu'à un objet appelable dont le déplacement est \texttt{noexcept}~? Quelles allocations restent-elles~? Qu'a \texttt{std::function} que votre classe n'a pas~?
\end{question}

\begin{question}[$\star\star\star\star$ \iqroles{QD, Dev, LLE} \iqfirm{toutes} \iqlevel{confirmé}]\label{iq:iv3:cpp:59}
En C++20, implémentez \texttt{Var<Ts.{}.{}.>}, un \texttt{std::variant} simplifié, avec \texttt{emplace<T>}, \texttt{get\_if<T>}, \texttt{index}, \texttt{visit} et le destructeur~; la copie peut être supprimée. Quelle taille le stockage doit-il avoir, et comment doit-il être aligné, pour une liste de types quelconque~? Comment le destructeur sait-il quoi détruire~? Dans quel état vous trouvez-vous si un constructeur lève une exception dans \texttt{emplace}~? Comment trouvez-vous l'indice d'un type à la compilation, sans récursion~?
\end{question}

\begin{question}[$\star\star\star\star$ \iqroles{QD, Dev, LLE, QR} \iqfirm{toutes} \iqlevel{confirmé}]\label{iq:iv3:cpp:60}
\texttt{d = a + b + c} sur des vecteurs de 1000 doubles, avec un \texttt{operator+} ordinaire renvoyant un \texttt{std::vector<double>}, alloue deux fois et exécute deux boucles d'addition. En C++20, écrivez des modèles d'expressions pour qu'il n'alloue rien, sachant que \texttt{d} a déjà sa taille, et qu'il n'exécute qu'une seule boucle. Que renvoie maintenant \texttt{a + b}~? Est-il sûr d'écrire \texttt{a = a + b}~? Maintenant, \texttt{auto e = a + b + c; use(e[0]);}~: que se passe-t-il, et pourquoi \texttt{auto} en est-il la cause~? Comment corrigeriez-vous la classe~?
\end{question}
